\section{Related Literature}\label{sec:literature}

Our result combines ingredients that are individually familiar. The contribution is therefore best located by separating the known components from the qualitative comparison that remains.

\paragraph{Strategic commitment and downstream equilibrium responses.}
\citet{LopezVives2019} study cost-reducing R\&D with spillovers and overlapping ownership. Their benchmark model is simultaneous, while a robustness extension lets firms commit to R\&D before downstream product-market competition. That two-stage extension illustrates a broad point relevant here: a first-stage investment incentive can depend on how the downstream equilibrium responds to investment. It does not contain a shared private fee setter or two decentralized public investors, so it is not a direct predecessor of our matched-price repricing comparison. We use it only as evidence that downstream-equilibrium feedback into investment incentives is a familiar strategic mechanism.

\paragraph{Public--private competition, pricing, and investment.}
\citet{Kim2024MixedDuopoly} is a particularly close institutional and strategic benchmark: a welfare-maximizing public platform competes with a profit-maximizing private platform in a two-sided linear-city model, and the paper compares first-stage R\&D incentives followed by platform price competition. It therefore establishes directly that mixed public--private platform competition, endogenous prices, network feedback, and investment incentives can coexist in one model. Its strategic object differs from ours, however. Kim compares the investment incentives of one public and one private platform; it does not study two decentralized public investors sharing a third private alternative, nor a matched-price experiment that switches the shared private price policy off and on while holding the public--public strategic object fixed.

\citet{ZogheibBourreau2021} study competition between profit-maximizing private and welfare-oriented public firms in network infrastructure, including local public firms and investment choices. \citet{Mun2019} and \citet{MunNakagawa2010} analyze regional or cross-border infrastructure provision with government pricing and investment. More recently, \citet{LiuReshidiRivadeneyra2026} study a welfare-maximizing public payment platform competing with a profit-maximizing private platform in a two-sided market and show that the private platform can adjust fees in response to public-platform competition. These papers make clear that mixed objectives, decentralized public investment, and private price responses to public competition are not themselves novel. Relative to this line, the present paper isolates how a shared private price policy changes the sign of the strategic relationship between two public investors.

\paragraph{Two-sided platforms, network interaction, and investment.}
Strategic interaction in two-sided market oligopoly is developed in \citet{FarhiHagiu2008}. \citet{HahnKim2026} study platform investment incentives under price-parity clauses and cross-platform investment externalities, while \citet{HeEtAl2026} analyze pricing and service investment under mixed network externalities and platform competition. These contributions absorb important components of our environment: cross-side network feedback, platform pricing, and strategic investment. We do not claim novelty for those ingredients.

\paragraph{Sequential pricing and qualitative reversals.}
\citet{BontemsHamiltonLepore2025} analyze sequential pricing on multisided platforms and show that timing changes equilibrium pricing and the role of cross-platform network effects. \citet{SanchezCartas2026} shows more sharply that a follower's best response in sequential two-sided platform competition can overturn the price-skewness ranking of the simultaneous benchmark. These papers make clear that a follower response can change qualitative conclusions in platform pricing. Our narrower object is different: the follower is a shared private intermediary, the upstream strategic actors are two welfare-oriented regional governments, and the comparison holds the private fee level fixed at the full-game on-path price before switching the private price policy on. Among the closest models examined in the prior-art audit, including the mixed-platform benchmark of \citet{Kim2024MixedDuopoly}, we do not identify a result that supplies this whole matched-price public--public best-response sign comparison by direct relabeling or restriction. This positioning is intentionally narrow and does not imply a new general theory of public platforms, sequential pricing, or two-sided markets.
